\documentclass[10pt,a4paper]{amsart}
\usepackage{amsmath,amssymb,amsthm,graphicx}
% For commutative diagrams we use Xy-pic
\usepackage[all]{xy}

% We use 2cell for 2-commutative diagrams.
\xyoption{2cell}
% Theorem environments.
%
\theoremstyle{plain}
\newtheorem{theorem}[subsection]{Theorem}
\newtheorem{proposition}[subsection]{Proposition}
\newtheorem{lemma}[subsection]{Lemma}

\theoremstyle{definition}
\newtheorem{definition}[subsection]{Definition}
\newtheorem{example}[subsection]{Example}
\newtheorem{exercise}[subsection]{Exercise}
\newtheorem{situation}[subsection]{Situation}

\theoremstyle{remark}
\newtheorem{remark}[subsection]{Remark}
\newtheorem{remarks}[subsection]{Remarks}

\numberwithin{equation}{subsection}

% Macros
%
\def\lim{\mathop{\mathrm{lim}}\nolimits}
\def\colim{\mathop{\mathrm{colim}}\nolimits}
\def\Spec{\mathop{\mathrm{Spec}}}
\def\Hom{\mathop{\mathrm{Hom}}\nolimits}
\def\Ext{\mathop{\mathrm{Ext}}\nolimits}
\def\SheafHom{\mathop{\mathcal{H}\!\mathit{om}}\nolimits}
\def\SheafExt{\mathop{\mathcal{E}\!\mathit{xt}}\nolimits}
\def\Sch{\mathit{Sch}}
\def\Mor{\mathop{\mathrm{Mor}}\nolimits}
\def\Ob{\mathop{\mathrm{Ob}}\nolimits}
\def\Sh{\mathop{\mathit{Sh}}\nolimits}
\def\NL{\mathop{N\!L}\nolimits}
\def\CH{\mathop{\mathrm{CH}}\nolimits}
\def\proetale{{pro\text{-}\acute{e}tale}}
\def\etale{{\acute{e}tale}}
\def\QCoh{\mathit{QCoh}}
\def\Ker{\mathop{\mathrm{Ker}}}
\def\Im{\mathop{\mathrm{Im}}}
\def\Coker{\mathop{\mathrm{Coker}}}
\def\Coim{\mathop{\mathrm{Coim}}}

% Boxtimes
%
\DeclareMathSymbol{\boxtimes}{\mathbin}{AMSa}{"02}

%
% Macros for moduli stacks/spaces
%
\def\QCohstack{\mathcal{QC}\!\mathit{oh}}
\def\Cohstack{\mathcal{C}\!\mathit{oh}}
\def\Spacesstack{\mathcal{S}\!\mathit{paces}}
\def\Quotfunctor{\mathrm{Quot}}
\def\Hilbfunctor{\mathrm{Hilb}}
\def\Curvesstack{\mathcal{C}\!\mathit{urves}}
\def\Polarizedstack{\mathcal{P}\!\mathit{olarized}}
\def\Complexesstack{\mathcal{C}\!\mathit{omplexes}}
% \Pic is the operator that assigns to X its picard group, usage \Pic(X)
% \Picardstack_{X/B} denotes the Picard stack of X over B
% \Picardfunctor_{X/B} denotes the Picard functor of X over B
\def\Pic{\mathop{\mathrm{Pic}}\nolimits}
\def\Picardstack{\mathcal{P}\!\mathit{ic}}
\def\Picardfunctor{\mathrm{Pic}}
\def\Deformationcategory{\mathcal{D}\!\mathit{ef}}
\usepackage{cleveref}
\crefname{originallemma}{original lemma}{original lemmas}
\Crefname{originallemma}{Original lemma}{Original lemmas}
\makeatletter
\newlabel{coherent-lemma-extend-coherent}{{0FD0}{source}{}{sourcepointer-1}{}}
\newlabel{coherent-lemma-extend-coherent@cref}{{[originallemma][0][]0FD0}{[1][0][]source}{}{}{}}
\makeatother
\crefname{originallemma}{original lemma}{original lemmas}
\Crefname{originallemma}{Original lemma}{Original lemmas}
\makeatletter
\newlabel{homology-lemma-serre-subcategory-K-groups}{{02MX}{source}{}{sourcepointer-2}{}}
\newlabel{homology-lemma-serre-subcategory-K-groups@cref}{{[originallemma][0][]02MX}{[1][0][]source}{}{}{}}
\makeatother
\crefname{originallemma}{original lemma}{original lemmas}
\Crefname{originallemma}{Original lemma}{Original lemmas}
\makeatletter
\newlabel{lemma-additivity-sheaf-cycle}{{02QZ}{source}{}{sourcepointer-3}{}}
\newlabel{lemma-additivity-sheaf-cycle@cref}{{[originallemma][0][]02QZ}{[1][0][]source}{}{}{}}
\makeatother
\crefname{originallemma}{original lemma}{original lemmas}
\Crefname{originallemma}{Original lemma}{Original lemmas}
\makeatletter
\newlabel{situation-setup}{{02QL}{source}{}{sourcepointer-4}{}}
\newlabel{situation-setup@cref}{{[originallemma][0][]02QL}{[1][0][]source}{}{}{}}
\makeatother
\title{Mutually inverse isomorphisms between the group of k-cycles and K\_0 of the truncated coherent category}
\author{The Stacks Project Authors}
\date{stacks-project:chow.tex:local-snapshot-20261012; GNU FDL 1.2 or later}
\begin{document}
\maketitle
\noindent\textit{Editorial note.} All delivered mathematical text is the author's own, copied from the pinned Stacks Project chapter file chow.tex; no statement or proof step was rewritten, reordered or removed.\par\medskip
\noindent\textit{Editorial note.} Preamble provenance: the pinned chapter file opens with \textbackslash{}input\{preamble\} and carries no preamble text of its own, so the project's single Stacks-wide preamble staged at sources/stacks-chow/preamble.tex is excerpted verbatim as two spans, lines 14-18 (Xy-pic setup) and lines 154-217 (the native theorem declarations and mathematical macros). Nothing was added, removed or redefined; the wrapper class and the amsmath/amssymb/amsthm/graphicx base packages are supplied by the assembler.\par\medskip
\noindent\textit{Editorial note.} Same-chapter pointers are located at their \textbackslash{}label lines inside the pinned chapter file; cross-chapter pointers are located in the project's own tag list staged in the snapshot (sources/stacks-chow/tags.txt), at the verbatim lines carrying their canonical tags.\par\medskip
\noindent\textit{Allowed standard prerequisites (editorial).}\begin{itemize}
\item Additivity of the cycle class of a coherent sheaf. Same chapter chow.tex, label lemma-additivity-sheaf-cycle, invoked at proof line 3790. Shows F -> [F]\_k is additive on Coh\_\{<=k\}(X), making the second map descend to K\_0.
\item Additive functions on an abelian category and K\_0 of a quotient category. Homology, label homology-lemma-serre-subcategory-K-groups part (1), invoked at proof line 3792. Supplies the formal criterion that an additive function killing Coh\_\{<=k-1\}(X) factors through K\_0 of the quotient; the actual factorisations are built in the counted span.
\item Extending a coherent sheaf from an open neighbourhood of a point. Cohomology of Schemes, label coherent-lemma-extend-coherent, invoked at proof line 3867. Produces the ideal sheaf I\_\{i,p\} inside O\_\{Z\_i\} mapping isomorphically onto F\textasciicircum{}p/F\textasciicircum{}\{p+1\} near the generic point xi\_i.
\end{itemize}\medskip
\ifdefined\hypertarget\hypertarget{sourcepointer-1}{}\fi
\noindent\textit{Source reference.} Original lemma 0FD0: sources/stacks-chow/tags.txt, source lines 19200--19200. Navigation pointer for \textbackslash{}ref\{coherent-lemma-extend-coherent\} used inside the counted span (line 3867); the argument that consumes it is retained in the span.\par\medskip
\ifdefined\hypertarget\hypertarget{sourcepointer-2}{}\fi
\noindent\textit{Source reference.} Original lemma 02MX: sources/stacks-chow/tags.txt, source lines 3460--3460. Navigation pointer for \textbackslash{}ref\{homology-lemma-serre-subcategory-K-groups\} used inside the counted span (line 3792); the argument that consumes it is retained in the span.\par\medskip
\ifdefined\hypertarget\hypertarget{sourcepointer-3}{}\fi
\noindent\textit{Source reference.} Original lemma 02QZ: sources/stacks-chow/pinned-original-source.tex, source lines 1789--1789. Navigation pointer for \textbackslash{}ref\{lemma-additivity-sheaf-cycle\} used inside the counted span (line 3790); the argument that consumes it is retained in the span.\par\medskip
\ifdefined\hypertarget\hypertarget{sourcepointer-4}{}\fi
\noindent\textit{Source reference.} Original lemma 02QL: sources/stacks-chow/pinned-original-source.tex, source lines 1372--1372. Navigation pointer for \textbackslash{}ref\{situation-setup\} used inside the counted span (line 3758); the argument that consumes it is retained in the span.\par\medskip
\begin{lemma}
\label{lemma-cycles-k-group}
Let $(S, \delta)$ be as in Situation \ref{situation-setup}.
Let $X$ be a scheme locally of finite type over $S$.
The maps
$$
Z_k(X)
\longrightarrow
K_0(\textit{Coh}_{\leq k}(X)/\textit{Coh}_{\leq k - 1}(X)),
\quad
\sum n_Z[Z] \mapsto
\left[\bigoplus\nolimits_{n_Z > 0} \mathcal{O}_Z^{\oplus n_Z}\right]
-
\left[\bigoplus\nolimits_{n_Z < 0} \mathcal{O}_Z^{\oplus -n_Z}\right]
$$
and
$$
K_0(\textit{Coh}_{\leq k}(X)/\textit{Coh}_{\leq k - 1}(X))
\longrightarrow
Z_k(X),\quad
\mathcal{F} \longmapsto [\mathcal{F}]_k
$$
are mutually inverse isomorphisms.
\end{lemma}

\begin{proof}
Note that if $\sum n_Z[Z]$ is in $Z_k(X)$, then
the direct sums
$\bigoplus\nolimits_{n_Z > 0} \mathcal{O}_Z^{\oplus n_Z}$ and
$\bigoplus\nolimits_{n_Z < 0} \mathcal{O}_Z^{\oplus -n_Z}$
are coherent sheaves on $X$ since the family $\{Z \mid n_Z > 0\}$
is locally finite on $X$.
The map $\mathcal{F} \to [\mathcal{F}]_k$ is additive
on $\textit{Coh}_{\leq k}(X)$, see
Lemma \ref{lemma-additivity-sheaf-cycle}. And $[\mathcal{F}]_k = 0$
if $\mathcal{F} \in \textit{Coh}_{\leq k - 1}(X)$. By part (1)
of Homology, Lemma \ref{homology-lemma-serre-subcategory-K-groups}
this implies that the second map is well defined too.
It is clear that the composition of the first map with the second
map is the identity.

\medskip\noindent
Conversely, say we start with a coherent sheaf $\mathcal{F}$
on $X$. Write $[\mathcal{F}]_k = \sum_{i \in I} n_i[Z_i]$
with $n_i > 0$ and $Z_i \subset X$, $i \in I$
pairwise distinct integral closed subschemes of $\delta$-dimension $k$.
We have to show that
$$
[\mathcal{F}] = [\bigoplus\nolimits_{i \in I} \mathcal{O}_{Z_i}^{\oplus n_i}]
$$
in $K_0(\textit{Coh}_{\leq k}(X)/\textit{Coh}_{\leq k - 1}(X))$.
Denote $\xi_i \in Z_i$ the generic point.
If we set
$$
\mathcal{F}' = \Ker(\mathcal{F} \to \bigoplus \xi_{i, *}\mathcal{F}_{\xi_i})
$$
then $\mathcal{F}'$ is the maximal coherent submodule of $\mathcal{F}$
whose support has dimension $\leq k - 1$. In particular $\mathcal{F}$
and $\mathcal{F}/\mathcal{F}'$ have the same class in
$K_0(\textit{Coh}_{\leq k}(X)/\textit{Coh}_{\leq k - 1}(X))$.
Thus after replacing $\mathcal{F}$ by $\mathcal{F}/\mathcal{F}'$
we may and do assume that the kernel $\mathcal{F}'$ displayed
above is zero.

\medskip\noindent
For each $i \in I$ we choose a filtration
$$
\mathcal{F}_{\xi_i} = \mathcal{F}_i^0 \supset \mathcal{F}_i^1 \supset
\ldots \supset \mathcal{F}_i^{n_i} = 0
$$
such that the successive quotients are of dimension $1$ over the residue
field at $\xi_i$. This is possible as the length of $\mathcal{F}_{\xi_i}$
over $\mathcal{O}_{X, \xi_i}$ is $n_i$.
For $p > n_i$ set $\mathcal{F}_i^p = 0$. For $p \geq 0$ we denote
$$
\mathcal{F}^p =
\Ker\left(\mathcal{F} \longrightarrow \bigoplus
\xi_{i, *}(\mathcal{F}_{\xi_i}/\mathcal{F}_i^p)\right)
$$
Then $\mathcal{F}^p$ is coherent, $\mathcal{F}^0 = \mathcal{F}$, and
$\mathcal{F}^p/\mathcal{F}^{p + 1}$ is isomorphic to a free
$\mathcal{O}_{Z_i}$-module of rank $1$ (if $n_i > p$) or $0$
(if $n_i \leq p$) in an open neighbourhood of $\xi_i$. Moreover,
$\mathcal{F}' = \bigcap \mathcal{F}^p = 0$. Since every quasi-compact
open $U \subset X$ contains only a finite number of $\xi_i$
we conclude that $\mathcal{F}^p|_U$ is zero for $p \gg 0$.
Hence $\bigoplus_{p \geq 0} \mathcal{F}^p$ is a coherent
$\mathcal{O}_X$-module. Consider the short exact sequences
$$
0 \to
\bigoplus\nolimits_{p > 0} \mathcal{F}^p \to
\bigoplus\nolimits_{p \geq 0} \mathcal{F}^p \to
\bigoplus\nolimits_{p > 0} \mathcal{F}^p/\mathcal{F}^{p + 1} \to 0
$$
and
$$
0 \to
\bigoplus\nolimits_{p > 0} \mathcal{F}^p \to
\bigoplus\nolimits_{p \geq 0} \mathcal{F}^p \to
\mathcal{F} \to 0
$$
of coherent $\mathcal{O}_X$-modules. This already shows that
$$
[\mathcal{F}] = [\bigoplus \mathcal{F}^p/\mathcal{F}^{p + 1}]
$$
in $K_0(\textit{Coh}_{\leq k}(X)/\textit{Coh}_{\leq k - 1}(X))$.
Next, for every $p \geq 0$ and $i \in I$ such that $n_i > p$
we choose a nonzero ideal sheaf $\mathcal{I}_{i, p} \subset \mathcal{O}_{Z_i}$
and a map $\mathcal{I}_{i, p} \to \mathcal{F}^p/\mathcal{F}^{p + 1}$ on $X$
which is an isomorphism over the open neighbourhood of $\xi_i$
mentioned above. This is possible by
Cohomology of Schemes, Lemma \ref{coherent-lemma-extend-coherent}.
Then we consider the short exact sequence
$$
0 \to
\bigoplus\nolimits_{p \geq 0, i \in I, n_i > p} \mathcal{I}_{i, p}
\to
\bigoplus \mathcal{F}^p/\mathcal{F}^{p + 1} \to
\mathcal{Q} \to 0
$$
and the short exact sequence
$$
0 \to
\bigoplus\nolimits_{p \geq 0, i \in I, n_i > p} \mathcal{I}_{i, p}
\to
\bigoplus\nolimits_{p \geq 0, i \in I, n_i > p} \mathcal{O}_{Z_i}
\to
\mathcal{Q}' \to 0
$$
Observe that both $\mathcal{Q}$ and $\mathcal{Q}'$ are zero in a neighbourhood
of the points $\xi_i$ and that they are supported on $\bigcup Z_i$.
Hence $\mathcal{Q}$ and $\mathcal{Q}'$ are in
$\textit{Coh}_{\leq k - 1}(X)$.
Since
$$
\bigoplus\nolimits_{i \in I} \mathcal{O}_{Z_i}^{\oplus n_i} \cong
\bigoplus\nolimits_{p \geq 0, i \in I, n_i > p} \mathcal{O}_{Z_i}
$$
this concludes the proof.
\end{proof}

\end{document}
